% IEEE-format paper (IEEEtran, conference layout with IEEE author blocks). Compile on Overleaf with pdfLaTeX. Upload the WHOLE folder
% (main.tex, references.bib, figures/architecture.png), e.g. the zip deliverables/paper/ieee_overleaf.zip.
% Status (1 Nov 2026): all experiments are final (Parts A-G: seven generators, TAT-QA, FinanceBench diagnosis, LLM judges, BioASQ). The only
% open item is the authors' verification of all 100 annotation labels (Section V-J, Human Validation); a ready replacement sentence is in a
% comment there.
% Caption rule (IEEE Editorial Style Manual): TABLE captions ABOVE the table (\caption before the tabular), figure captions BELOW the figure.
\documentclass[conference]{IEEEtran}
\usepackage[T1]{fontenc}
\usepackage{cite}
\usepackage{amsmath,amssymb}
\usepackage{graphicx}
\usepackage{booktabs}
\usepackage{multirow}
\usepackage{array}
\usepackage{xcolor}
\usepackage{url}
\usepackage[hidelinks]{hyperref}

\newcommand{\pending}[1]{\textcolor{red}{[PENDING: #1]}}
\newcommand{\m}[1]{\texttt{#1}}   % method / signal names

\begin{document}

\title{Cross-Domain Attribution, Hallucination Detection and Refusal in Retrieval-Augmented Generation}

% Authors in the IEEE author-block style. Since all authors share the same affiliation, list it once at the bottom.
\author{%
\IEEEauthorblockN{Utkarsh Kumar Jha, Surveer Singh Guglani, Prajjwal Joshi, and Dr. Savita Yadav}
\IEEEauthorblockA{Dept. of Computer Science and Engineering\\
Netaji Subhas University of Technology\\
New Delhi, India\\
\{utkarsh.kumar\_ug23, surveer.singh.ug23, prajjwal.joshi.ug23, savita.yadav\}@nsut.ac.in}
}

\maketitle

\begin{abstract}
Retrieval-augmented generation (RAG) is increasingly used on medical, legal and financial documents, but the mechanisms that make it trustworthy,
namely evidence attribution, hallucination detection and refusal, have mostly been built and tested on general-knowledge text. We test whether these
mechanisms keep working when only the domain changes. We build a testbed of 1,350 questions over five corpora (HotpotQA, SQuAD~2.0, PubMedQA,
CUAD and the full FinanceBench filings) in which the gold evidence chunk is known for 999 of 1,000 answerable questions, every generator sees the
same retrieved passages, and each question is also posed with the gold passage guaranteed (oracle) and with it removed. Across seven open-weight
generators from five families (1.5B to 32B parameters), we find that: (i) retrieval is the largest source of error in legal and financial text, while
medical errors come from the generator; (ii) even with the gold passage guaranteed, attribution accuracy falls from about 0.83--0.93 on general text to
0.32--0.39 on financial filings for every model, and the best attribution method reverses by domain; (iii) the collapse replicates on a second
financial dataset (TAT-QA, 7 of 7 models, pre-registered) and has an identifiable cause: similarity fails on table cells, and the decisive signal is
the numbers. A label-free, model-free number-aware scorer raises financial attribution by 0.12--0.24, confirmed on held-out TAT-QA data (0.51 to 0.75)
without changing the other domains. LLM judges (8B and 32B, equally good) remain more accurate (0.69--0.79 on financial text), and the free scorer
closes 67--87\% of the gap to them; (iv) on a second medical dataset with sentence-level evidence (BioASQ), attribution is 0.62--0.71, well below
what PubMedQA's whole-abstract gold suggests, for all seven models; (v) refusing when evidence is missing improves with model size and transfers
across four domains, but over-refusal does not decrease; (vi) token log-probabilities are the best single signal for answering only the most confident half of questions, but do not become
more informative with size, and no signal beats trivial policies when wrong answers are costly; and (vii) a supervised hallucination checker that
improves on our zero-shot detector on RAGTruth (51.4 vs.\ 43.6 F1, $p=0.001$) does not transfer to our domains. We use pre-registered,
multiplicity-controlled tests and a measured re-run noise floor.
\end{abstract}

\begin{IEEEkeywords}
Retrieval-augmented generation, evidence attribution, hallucination detection, selective prediction, domain shift, large language models.
\end{IEEEkeywords}

% =====================================================================================================================
\section{Introduction}
Large language models (LLMs) write fluent answers that are not always correct \cite{ji2023survey,huang2025survey}.
Retrieval-augmented generation (RAG) \cite{lewis2020rag} reduces this problem by first retrieving passages from a document collection and then asking
the model to answer from them. RAG is now a common design for assistants in healthcare, law and finance. Retrieval alone, however, does not make the
answers trustworthy. Commercial legal research tools that use retrieval still give incorrect or unsupported answers to 17--33\% of queries
\cite{magesh2025legal}, and general-purpose models invent legal facts even more often \cite{dahl2024legal}. A trustworthy RAG system therefore needs
three more abilities: it should show which passage supports each part of its answer (\emph{attribution}), flag statements that no passage supports
(\emph{hallucination detection}), and decline to answer when the evidence is not enough (\emph{refusal}).

Most methods for these three tasks were developed and evaluated on Wikipedia-style questions
\cite{gao2023alce,song2025trust,niu2024ragtruth,joren2025sufficient}. Specialised documents are different. Contracts are long and formulaic,
biomedical abstracts require the reader to draw a conclusion, and annual reports are full of tables and numbers. It is not known whether methods
tuned on general text still work on such documents. Existing domain-specific studies are hard to compare with each other because each one uses its
own retriever, prompts, model and metrics.

In this work we hold the whole pipeline fixed and change only the domain. Our contributions are:
\begin{enumerate}
\item A \textbf{controlled cross-domain testbed} over general, medical, legal and financial text with gold-evidence chunk labels, retrieval computed
once so that every model sees identical passages, and two counterfactual conditions: an \emph{oracle} condition that guarantees the gold passage and
an \emph{evidence-removal} condition that deletes it. Together they separate retrieval failures from generation and attribution failures and give
ground-truth unanswerable questions in every domain.
\item A \textbf{retrieval-free measurement of the attribution gap}. With the gold passage guaranteed, attribution accuracy on financial filings stays
near one third for all seven models, and the best method changes with the domain. A per-domain mixture chosen from 20 labelled questions repairs the
worst case and transfers across generators.
\item A \textbf{diagnosis and a fix}. Controlled re-scoring of the same answers shows that table cells, not computation or boilerplate, break
similarity-based attribution, and that the numbers in a sentence are the missing signal. A number-aware scorer that needs no labels and no model,
designed on FinanceBench, is confirmed on a second financial dataset (TAT-QA) under hypotheses registered before that data was generated, and is
compared with LLM judges: it closes most, but not all, of the gap to them at no cost.
\item \textbf{Replication on second datasets} for the two domains where the first dataset was a weak test: TAT-QA (finance) and BioASQ with
sentence-level expert evidence (medicine), each with hypotheses registered in advance.
\item A \textbf{model-scale study} with seven open-weight generators (1.5B--32B, five families) of refusal and confidence signals, tested against
hypotheses that were written down before the larger models were run.
\item \textbf{Negative results and method checks} that we believe are useful to others: a measured noise floor for repeated GPU runs, a supervised
hallucination checker that does not transfer, and an end-to-end ``trust layer'' that does not beat trivial policies once it is evaluated fairly.
\end{enumerate}
Code, configurations, the pre-registered analysis plans and all aggregate results are released at
\url{https://github.com/Surveerssg/trustworthy-rag-attribution}.

% =====================================================================================================================
\section{Related Work}
\label{sec:related}
Table~\ref{tab:related} summarises the closest prior work and what each leaves open for our question.

\begin{table*}[!t]
\caption{Closest Related Work and the Dimensions Covered by This Study}
\label{tab:related}
\centering
\footnotesize
\begin{tabular}{@{}p{3.4cm}p{3.9cm}p{3.1cm}p{6.0cm}@{}}
\toprule
Work & Contribution & Evaluated on & Gap relative to this study \\
\midrule
RAG \cite{lewis2020rag} & Retrieve-then-generate architecture & Open-domain QA & No attribution, detection or refusal \\
AIS \cite{rashkin2023ais} & Definition of attributable statements & General NLG & Human judgement; not domain-controlled \\
ALCE \cite{gao2023alce} & Citation recall/precision with NLI & General QA & General domain only; no gold evidence chunk \\
MIRAGE \cite{qi2024mirage} & Attribution from model internals & Extractive / open QA & White-box access; general domain \\
Wallat et al.\ \cite{wallat2025correctness} & Correct citations are not always faithful & General QA & Single domain \\
Trust-Align \cite{song2025trust} & Training to cite and to refuse & ASQA, QAMPARI, ELI5 & Requires fine-tuning; general domain \\
RAGTruth \cite{niu2024ragtruth} & Human-labelled RAG hallucinations & QA, summaries, data-to-text & Detection only; general domain \\
RGB \cite{chen2024rgb} & Noise robustness and negative rejection & General QA (en/zh) & No attribution; no specialised domains \\
Sufficient Context \cite{joren2025sufficient} & Context sufficiency for selective answers & General QA & Single domain; no attribution \\
RAGBench \cite{friel2024ragbench}, RAG-QA Arena \cite{han2024ragqa} & Multi-domain RAG evaluation & Industry / long-form QA & No gold evidence chunk; no counterfactual contexts \\
MedRAG \cite{xiong2024medrag}, LegalBench-RAG \cite{pipitone2024legalbenchrag}, FinanceBench \cite{islam2023financebench} & Domain benchmarks & One domain each & Answer quality within one domain \\
\midrule
\textbf{This work} & Attribution, detection and refusal under one fixed pipeline & General, medical, legal, financial (two datasets each for medicine and finance) & Gold chunks, oracle and removal conditions, seven generators, a diagnosed fix compared with LLM judges, free compute \\
\bottomrule
\end{tabular}
\end{table*}

\subsection{Hallucination and Retrieval-Augmented Generation}
Surveys separate hallucinations caused by training data from those introduced while generating \cite{ji2023survey} and list open problems that are
specific to RAG \cite{huang2025survey}. Retrieval reduces unsupported output but does not remove it \cite{lewis2020rag,magesh2025legal}. Practitioners
report missing content, relevant documents ranked too low and wrong extraction as common failure points \cite{barnett2024seven}. Models also use
evidence unevenly depending on where it appears in the context \cite{liu2024lost}, and models can be trained to decide when to retrieve and to
critique their own output \cite{asai2024selfrag}.

\subsection{Evidence Attribution}
Attributable to Identified Sources (AIS) defines when a statement is supported by a source \cite{rashkin2023ais}. ALCE measures citation recall and
precision with an NLI model \cite{gao2023alce}, following earlier NLI-based factual-consistency evaluation \cite{honovich2022true}. Model-internal
methods attribute answers using attention and gradients \cite{qi2024mirage}, and perturbation methods such as leave-one-out and Shapley values measure
causal influence at a much higher cost \cite{nematov2025source}. Wallat et al.\ show that a correct citation is not necessarily the source the model
used \cite{wallat2025correctness}. All of these studies use general-knowledge data.

\subsection{Hallucination Detection}
Black-box detectors compare several sampled answers \cite{manakul2023selfcheck} or measure semantic entropy \cite{farquhar2024semantic}. RAG-specific
detectors check the answer against the retrieved passages with NLI models or LLM judges \cite{es2024ragas}. RAGTruth provides human span-level labels
for RAG outputs; on its QA split, prompting GPT-4-turbo reaches 45.6 F1 and a fine-tuned Llama-2-13B reaches 68.2 \cite{niu2024ragtruth}.

\subsection{Confidence, Selective Answering and Refusal}
Selective prediction trades coverage for risk \cite{geifman2017selective,geifman2019bias}, and conformal risk control chooses thresholds with
finite-sample guarantees \cite{angelopoulos2024crc}. Proposed refusal signals include the model's own knowledge of what it knows \cite{kadavath2022know},
whether the context is sufficient \cite{joren2025sufficient}, and training to refuse \cite{song2025trust}. An axiomatic analysis shows that common
uncertainty measures behave inconsistently when retrieval is added \cite{soudani2025uncertainty}. RGB tests whether models reject questions when the
evidence is missing \cite{chen2024rgb}.

\subsection{Domain-Specific RAG}
There are domain benchmarks for medicine \cite{xiong2024medrag}, law \cite{pipitone2024legalbenchrag,hendrycks2021cuad} and finance
\cite{islam2023financebench}, and multi-domain collections that score answer quality \cite{friel2024ragbench,han2024ragqa}. These measure end-task
quality within or across domains, but they do not hold the pipeline fixed, do not give gold evidence chunks for measuring attribution accuracy, and do
not separate retrieval failure from attribution failure. This combination is the gap we address.

% =====================================================================================================================
\section{Testbed}
\label{sec:testbed}

\subsection{Datasets and Gold Evidence}
Table~\ref{tab:data} lists the five datasets. We converted all questions to one format. For every answerable question we located the chunk that holds
the gold evidence using the character offset, an exact substring match or a token recall of at least 0.5. This succeeded for 999 of 1,000 answerable
questions; for PubMedQA every chunk of the source abstract counts as gold. PubMedQA abstracts exclude the conclusion, so the answer has to be inferred
\cite{jin2019pubmedqa}. CUAD questions name the contract because the corpus pools 60 contracts, and clause types that lawyers marked as absent give
unanswerable questions \cite{hendrycks2021cuad}. The FinanceBench corpus contains every page of all 84 filings, not only the evidence pages
\cite{islam2023financebench}, so retrieval is realistic.

\textbf{Second financial dataset.} To test whether the financial results are a property of FinanceBench, we added TAT-QA \cite{zhu2021tatqa}: hybrid
excerpts of annual reports, each a table with its surrounding paragraphs. All 278 development-set contexts form one pooled corpus (one document each:
the table linearised row by row, then its paragraphs; same chunker, embedder and prompts). TAT-QA labels every answer by its source (table, text, or
both) and by kind (computed by arithmetic or counting, or quoted as a span), and these labels give the gold evidence (the table, the relevant
paragraphs, or both). We sampled 200 questions stratified over source $\times$ kind (37 per cell; only 16 text-computed questions exist). TAT-QA
questions do not name the company, so natural retrieval over the pooled corpus is ill-posed; we use only the oracle and removal conditions.

\textbf{Second medical dataset.} PubMedQA counts the whole abstract as gold, which makes medical attribution easy to score. BioASQ
\cite{krithara2023bioasq} (Task B training set 13b) gives expert snippets, usually one or two sentences, inside PubMed abstracts. We sampled 200
questions (50 yes/no answered yes, 50 answered no, 100 factoid) whose snippets are found in the current abstract. The corpus holds the title and
abstract of every document listed for these questions plus those of 800 other BioASQ questions (12,837 abstracts, fetched from PubMed). Gold chunks are
the chunks that contain a snippet; the removal condition deletes every document listed for the question. BioASQ questions are self-contained, so all
conditions apply.

\begin{table}[!t]
\caption{Datasets, Corpora and Questions}
\label{tab:data}
\centering
\footnotesize
\begin{tabular}{@{}llp{2.1cm}p{1.8cm}@{}}
\toprule
Domain & Dataset & Corpus searched & Questions \\
\midrule
General & HotpotQA \cite{yang2018hotpotqa} & 9,771 Wikipedia paragraphs & 200 (multi-hop) \\
General & SQuAD 2.0 \cite{rajpurkar2018squad2} & 1,204 paragraphs & 500 (250 unanswerable) \\
Medical & PubMedQA \cite{jin2019pubmedqa} & 1,000 abstracts, conclusions removed & 200 (yes/no/maybe) \\
Medical & BioASQ \cite{krithara2023bioasq} & 12,837 PubMed abstracts & 200 (yes/no, factoid) \\
Legal & CUAD \cite{hendrycks2021cuad} & 60 full contracts & 300 (100 clause absent) \\
Financial & FinanceBench \cite{islam2023financebench} & 11,917 pages of 84 reports & 150 (numeric) \\
Financial & TAT-QA \cite{zhu2021tatqa} & 278 table + text excerpts & 200 (oracle, removal) \\
\bottomrule
\end{tabular}
\end{table}

\subsection{Pipeline}
Fig.~\ref{fig:arch} shows the pipeline. Documents are split at sentence boundaries into chunks of about 180 words with a one-sentence overlap (the
embedding model truncates at 256 word pieces). Chunks are embedded with all-MiniLM-L6-v2 \cite{reimers2019sbert,wang2020minilm} and indexed separately
for each domain. The top five chunks for every question and condition are computed once and stored, so every generator receives exactly the same
passages. We use two prompts that are identical in every domain except for a one-line answer-format hint: a \emph{vanilla} prompt (``use the following
passages'') and a \emph{trust} prompt that asks for a passage citation after every sentence and the reply ``INSUFFICIENT EVIDENCE'' when the passages
are not enough. Generation uses temperature 0, at most 200 new tokens and a 4,096-token context window.

We use seven open-weight instruction-tuned generators from five families, all served by Ollama with 4-bit weights on free Kaggle GPUs (2$\times$ NVIDIA
T4): Qwen2.5-1.5B, Qwen2.5-14B and Qwen2.5-32B \cite{qwen2024qwen25}, Llama-3.1-8B \cite{grattafiori2024llama3}, Phi-4-14B \cite{abdin2024phi4},
Mistral-Small-24B \cite{mistral2025small3} and Gemma-3-27B \cite{gemma2025gemma3}. The 1.5B, 8B and Qwen 14B models answer all 1,350 questions. To fit
Kaggle's 12-hour session limit, the other four answer a fixed random sample of 100 questions per domain (500 questions, the same sample for all four,
with the answerable/unanswerable ratio kept); the 24B--32B models have their layers split across both GPUs. Every comparison that involves these
models is also reported with all models restricted to this shared subset, so the comparison is paired. All seven models answer the 200 TAT-QA and the
200 BioASQ questions.

\begin{figure}[!t]
\centering
\includegraphics[width=\columnwidth]{figures/paratrust_flowchart.pdf}
\caption{Pipeline: offline chunking and indexing, retrieval computed once for every condition, generation, and the trust layer (attribution,
hallucination detection and confidence signals).}
\label{fig:arch}
\end{figure}

\subsection{Retrieval Conditions}
Each question is answered under four retrieval conditions:
(i) \emph{base}: dense top-5 retrieval;
(ii) \emph{improved}: a retriever chosen per domain: hybrid BM25 \cite{robertson2009bm25} and dense retrieval combined with reciprocal rank fusion
\cite{cormack2009rrf} and cross-encoder re-ranking \cite{nogueira2019rerank} for general text, plain dense retrieval for medical text, and hybrid
retrieval restricted to the named contract or filing for legal and financial text;
(iii) \emph{oracle}: the gold chunk is placed in the top five at a random position;
(iv) \emph{removal}: every gold chunk (for PubMedQA, the whole source abstract) is removed from the candidates, so the correct behaviour is to refuse.
Removal is applied only to answerable questions.

\subsection{Correctness and Attribution Labels}
An answer is correct if it contains the gold answer as a whole word (short answers), starts with the gold label (PubMedQA), recalls at least half of the
gold clause's tokens (CUAD), or contains the gold number within 1\% (FinanceBench; otherwise an NLI check is used). Refusals are explicit
(``INSUFFICIENT EVIDENCE'', ``I don't know'') or implicit (the first sentence says the passages do not contain the answer; all 25 detections we checked
by hand were correct). \emph{Attribution accuracy} is the share of answer sentences whose top-scored chunk is a gold chunk. It is computed only when a
gold chunk is in the context.

% =====================================================================================================================
\section{Methods}
\label{sec:methods}

\subsection{Attribution Methods}
Every answer sentence is scored against every retrieved chunk. \m{sim\_chunk} and \m{sim\_sent} are the cosine similarity with the whole chunk and with
its most similar sentence. \m{nli\_chunk} is the entailment probability of a DeBERTa-v3 NLI cross-encoder \cite{he2023debertav3} with the chunk as the
premise. We added two variants for long, dense text: \m{nli\_window} takes the highest entailment over two-sentence windows, and \m{nli\_set} also tests
the sentence against the two most similar chunks joined together, so that a claim supported only by two passages together is still recognised. In
\m{nli\_set} both chunks of the pair receive the joint score, and ties are broken by each chunk's own single-chunk score
(Section~\ref{sec:neg}). Leave-one-out (LOO) removes one chunk, generates again and checks whether the answer changes; we use it as a causal reference.
A number-aware scorer that we derived from the FinanceBench failures is described with its diagnosis in Section~\ref{sec:why}.

\subsection{Per-Domain Attribution Mixture}
Since no single method is best in every domain (Section~\ref{sec:attr}), we score each chunk as
$s = w\,\m{sim\_sent} + (1-w)\,\m{nli\_set}$ with $w\in\{0,0.1,\dots,1\}$ chosen for each domain from 20 labelled questions (ties go to the larger
$w$, which needs no NLI model). Weights fitted on Llama-3.1-8B oracle questions are 1.0 for SQuAD, HotpotQA and PubMedQA, 0.7 for CUAD and 0.0 for
FinanceBench. The same weights are applied unchanged to the other four generators. Section~\ref{sec:attr} also reports how the choice depends on
the number of labelled questions.

\subsection{Confidence Signals}
For each answer we compute retrieval signals (top-1 similarity \m{ret\_top1}, mean, gap), NLI coverage (the mean and minimum over sentences of the best
support), similarity coverage, token log-probability summaries (mean, first tokens, minimum, share of low-probability tokens, top-2 margin), and the score
of a supervised hallucination checker (Section~\ref{sec:ragtruth}). The primary log-probability signal, fixed in advance, is \m{lp\_mean}, the mean token
log-probability of the answer.

\subsection{Evaluation Protocol}
\label{sec:protocol}
Confidence signals are compared with two trivial policies, answering everything and refusing everything, using a per-question utility of $+1$ for a
correct answer or a correct refusal, $0$ for refusing an answerable question and $-\lambda$ for a wrong answer, with $\lambda\in\{0.5,1,2,4\}$. Thresholds
are chosen with repeated five-fold cross-fitting, so each question is decided by a threshold fitted on other questions. Differences are tested with a
paired bootstrap over questions (10,000 resamples) \cite{efron1993bootstrap}. Families of tests are controlled with the Benjamini--Hochberg procedure
at a 10\% false discovery rate \cite{benjamini1995fdr}, and every test is run against a deliberately strong control: the best
non-log-probability signal of each cell, picked on the same data.

Before the 14B, 27B and 32B models were run we recorded four hypotheses: \textbf{H1} log-probabilities become more informative with model size;
\textbf{H2} they become useful at high error cost ($\lambda\geq2$); \textbf{H3} they beat retrieval and NLI signals on HotpotQA and PubMedQA but not on
CUAD and FinanceBench; \textbf{H4} refusal improves with size while over-refusal does not grow. We also fixed that comparisons involving the 27B and 32B
models are repeated on the shared subset. Phi-4 and Mistral-Small, added later, were analysed with the same tests, unchanged. Before any TAT-QA answer was generated we recorded a second plan (T1--T7, Section~\ref{sec:tatqa}) with its
tests: cluster bootstrap over questions (10,000 resamples, all models' records of a question resampled together), Holm correction within each family,
and an effect is claimed only if it has the predicted sign, $|\Delta|\geq0.05$ and a corrected $p<0.05$. The same rules govern two later plans,
each written before its data existed: an LLM-as-attributor baseline (L1--L4, Section~\ref{sec:judge}) and BioASQ (B1--B5, Section~\ref{sec:bioasq}).

\textbf{Noise floor.} Running an identical configuration twice on the same hardware (same contexts, temperature 0, fixed seed) changed the answer text
for 38\% (Llama) and 45\% (Qwen-1.5B) of questions and changed correctness for 1.7\% and 3.3\%. In one domain, accuracy moved by 4.0 points with a
bootstrap $p=0.01$ although the inputs were identical. We therefore treat accuracy differences smaller than 4--5 points as noise whatever the bootstrap
says.

% =====================================================================================================================
\section{Results}
\label{sec:results}

\subsection{Retrieval and the Oracle Decomposition}
\label{sec:retrieval}
Retrieval quality differs by an order of magnitude across domains (Table~\ref{tab:retrieval}). The domain-specific retriever raises Hit@5 everywhere it
applies (FinanceBench 0.19 to 0.31, CUAD 0.59 to 0.71) and improves accuracy by 7--11 points in the cells where the confidence interval excludes zero
and the gain is above the noise floor. On questions where the gold chunk was retrieved by both retrievers, attribution accuracy changed by at most
$\pm0.07$ with no consistent direction, so better retrieval did not make attribution worse.

\begin{table}[!t]
\caption{Hit@5 and Accuracy Gain From Domain-Specific Retrieval (Vanilla Prompt, All Questions)}
\label{tab:retrieval}
\centering
\footnotesize
\setlength{\tabcolsep}{3.5pt}
\begin{tabular}{@{}lccccc@{}}
\toprule
& \multicolumn{2}{c}{Hit@5} & \multicolumn{3}{c}{$\Delta$Accuracy, base $\rightarrow$ improved} \\
\cmidrule(lr){2-3}\cmidrule(l){4-6}
Domain & base & improved & Qwen-1.5B & Llama-8B & Qwen-14B \\
\midrule
SQuAD 2.0 & .872 & .992 & +.072* & +.084* & +.092* \\
HotpotQA & .965 & 1.00 & +.070* & +.060 & +.070* \\
PubMedQA & .990 & .990 & -- & -- & -- \\
CUAD & .585 & .710 & +.045 & +.075* & +.105* \\
FinanceBench & .188 & .309 & +.013 & +.080* & +.067 \\
\bottomrule
\multicolumn{6}{@{}p{0.95\columnwidth}@{}}{\scriptsize * 95\% paired-bootstrap CI excludes zero. PubMedQA uses dense retrieval in both arms.}
\end{tabular}
\end{table}

The oracle condition shows how much accuracy is lost because of retrieval (Table~\ref{tab:oracle}). For every model of 8B or more, the loss is
+0.22 to +0.29 in FinanceBench and +0.21 to +0.30 in CUAD. In PubMedQA it is close to zero: the gold abstract is retrieved 99\% of the time and
guaranteeing it does not help, so the errors come from the model's inference, and accuracy stays at 0.52--0.57 even for the 24B--32B models. With the
base retriever, FinanceBench accuracy is 0.24--0.28 for all seven models: a larger model cannot make up for evidence that was never retrieved.

\begin{table*}[!t]
\caption{Accuracy With the Gold Passage Guaranteed (Oracle) and, in Brackets, the Loss Due to Retrieval (Oracle $-$ Base), Vanilla Prompt}
\label{tab:oracle}
\centering
\footnotesize
\setlength{\tabcolsep}{4pt}
\begin{tabular}{@{}lccccccc@{}}
\toprule
Domain & Qwen-1.5B & Llama-8B & Qwen-14B & Phi-4-14B$^\dagger$ & Mistral-24B$^\dagger$ & Gemma-27B$^\dagger$ & Qwen-32B$^\dagger$ \\
\midrule
SQuAD & .808 (.06) & .940 (.10) & .868 (.08) & .900 (.14) & .940 (.14) & .880 (.16) & .900 (.14) \\
HotpotQA & .605 (.10) & .765 (.16) & .795 (.17) & .800 (.19) & .810 (.17) & .790 (.18) & .780 (.16) \\
PubMedQA & .380 ($-$.09) & .540 ($-$.02) & .570 (.01) & .530 ($-$.04) & .570 (.03) & .520 ($-$.03) & .540 (.03) \\
CUAD & .610 (.14) & .680 (.21) & .645 (.21) & .776 (.25) & .701 (.30) & .672 (.30) & .776 (.28) \\
FinanceBench & .336 (.09) & .530 (.29) & .550 (.29) & .576 (.29) & .475 (.22) & .535 (.29) & .545 (.29) \\
\bottomrule
\multicolumn{8}{@{}p{0.95\textwidth}@{}}{\scriptsize $^\dagger$ 100 questions per domain (shared subset). Retrieval losses of the 8B model have 95\% CIs that
exclude zero in SQuAD, HotpotQA, CUAD and FinanceBench.}
\end{tabular}
\end{table*}

\subsection{Attribution Across Domains}
\label{sec:attr}
Table~\ref{tab:attr} gives attribution accuracy in the oracle condition, where the gold chunk is always present, so retrieval cannot explain the
differences. On general text, sentence similarity reaches 0.83--0.93. On financial filings it reaches only 0.32--0.39 for \emph{every} model (chance is
about 0.2), so the gap is neither a retrieval problem nor a model-size or model-family problem. The best method also depends on the domain: sentence
similarity is better on general and medical text, while evidence-set NLI is clearly better on financial filings (0.51--0.55). Section~\ref{sec:why}
identifies the cause. The 20-label mixture uses each method where it works. With the weights fitted once on Llama and applied unchanged, it raises the
worst domain from 0.32--0.39 to 0.50--0.55 and the macro average from 0.70--0.75 to 0.74--0.81 for all seven generators. For Llama, where the weights
were fitted, held-out estimates over 200 random draws of the 20 labelled questions give a macro score of 0.737 and a worst domain of 0.483. In
natural retrieval order the mixture performs about the same as \m{sim\_sent} on average (0.746 vs.\ 0.740 for Llama), because the gain is concentrated
in the domain where similarity fails.

\begin{table*}[!t]
\caption{Attribution Accuracy With the Gold Passage Guaranteed (Oracle Condition): \m{sim\_sent} / \m{nli\_set}, and the 20-Label Mixture}
\label{tab:attr}
\centering
\footnotesize
\setlength{\tabcolsep}{5pt}
\begin{tabular}{@{}lccccccc@{}}
\toprule
Domain & Qwen-1.5B & Llama-8B & Qwen-14B & Phi-4-14B$^\dagger$ & Mistral-24B$^\dagger$ & Gemma-27B$^\dagger$ & Qwen-32B$^\dagger$ \\
\midrule
SQuAD & .90/.87 & .87/.88 & .90/.88 & .83/.78 & .84/.84 & .84/.80 & .89/.85 \\
HotpotQA & .85/.78 & .88/.80 & .91/.81 & .89/.86 & .89/.84 & .89/.82 & .93/.86 \\
PubMedQA & .88/.63 & .74/.61 & .89/.64 & .80/.54 & .80/.64 & .69/.55 & .89/.60 \\
CUAD & .58/.60 & .71/.60 & .69/.65 & .69/.66 & .73/.57 & .74/.67 & .72/.63 \\
FinanceBench & \textbf{.32}/.53 & \textbf{.35}/.51 & \textbf{.36}/.54 & \textbf{.39}/.55 & \textbf{.38}/.51 & \textbf{.35}/.52 & \textbf{.32}/.55 \\
\midrule
Mixture, macro & .751 & .743 & .792 & .755 & .755 & .741 & .808 \\
Mixture, worst & .528 & .505 & .535 & .549 & .510 & .521 & .551 \\
\m{sim\_sent}, worst & .322 & .347 & .364 & .390 & .384 & .354 & .318 \\
\bottomrule
\multicolumn{8}{@{}p{0.95\textwidth}@{}}{\scriptsize $^\dagger$ 100 questions per domain. Mixture weights: 1.0 (SQuAD, HotpotQA, PubMedQA), 0.7 (CUAD),
0.0 (FinanceBench), fitted on 20 Llama questions per domain and applied unchanged to the other six generators.}
\end{tabular}
\end{table*}

\textbf{How many labels are needed?} We repeated the weight choice with 0, 2, 5, 10, 20 and 50 labelled questions (500 random draws per model
and domain; within a draw every label budget is scored on the same held-out questions, and 0 labels means the default $w=0.9$). Table~\ref{tab:labels}
shows that labels matter only where the domain behaves differently. On the four non-financial domains the default weight is already within 0.01 of
the best fixed weight, and small label budgets do not help (they occasionally pick a worse weight). On FinanceBench, 20 labels raise attribution
accuracy from 0.37 to 0.51, 87\% of the attainable gain (the ceiling, 0.53, equals pure \m{nli\_set}), and 50 labels reach the ceiling; with 5--10
labels an unlucky draw gains almost nothing. The mixture therefore does not beat the best single method within a domain: its value is choosing the
right method for each domain from a handful of labelled questions. This analysis is exploratory; the 20-label recipe was fixed before it was run.

\begin{table}[!t]
\caption{Held-Out Attribution Accuracy (Oracle Condition) vs.\ the Number of Labelled Questions Used to Choose the Mixture Weight}
\label{tab:labels}
\centering
\footnotesize
\setlength{\tabcolsep}{4pt}
\begin{tabular}{@{}lccc@{}}
\toprule
Labels & Other 4 domains & FinanceBench & FinanceBench, 10th pct. \\
\midrule
0 (default $w=0.9$) & .822 & .367 & .328 \\
2 & .808 & .431 & .331 \\
5 & .810 & .465 & .349 \\
10 & .812 & .489 & .379 \\
20 & .815 & .508 & .431 \\
50 & .819 & .525 & .486 \\
\midrule
Best fixed $w$ (ceiling) & .831 & .529 & -- \\
\m{sim\_sent} only & .813 & .355 & -- \\
\m{nli\_set} only & .729 & .529 & -- \\
\bottomrule
\multicolumn{4}{@{}p{0.95\columnwidth}@{}}{\scriptsize Mean over the seven generators and 500 draws. The four 100-question models are omitted on SQuAD
and CUAD (too few questions for a 50-question label pool). 10th pct.\ = an unlucky label draw.}
\end{tabular}
\end{table}

\textbf{Faithfulness.} For 40 answers per domain (Llama-8B and Qwen-1.5B) we ran leave-one-out. When a method points at the gold chunk, removing that
chunk changes the answer in 73--100\% of cases ($\tau=0.5$), so ``cited but not used'' is rare. Finding the \emph{most} influential chunk is harder
(0.35--0.60 on general and medical text for Llama, 0.18--0.55 on CUAD and 0.28--0.38 on FinanceBench). Of the methods we compared, \m{sim\_sent} agrees
best with the causal influence in 9 of 10 model--domain cells (Spearman 0.11--0.51).

\subsection{Why Attribution Fails on Financial Filings, and a Fix}
\label{sec:why}
We diagnosed the FinanceBench failures without regenerating anything: the 842 oracle answers of the seven models stay fixed, and only the passage
text given to the attribution scorer is changed, identically for gold and distractor passages. This analysis was designed after looking at
FinanceBench and is exploratory; its held-out test is TAT-QA (Section~\ref{sec:tatqa}).

\textbf{Diagnosis.} PDF text extraction puts every table cell on its own line, so the sentence splitter hands the similarity scorer fragments such as
``1,488'' or ``\$''. When at least 30\% of the gold passage's lines are bare numbers (600 of 842 answers), \m{sim\_sent} falls to 0.28, against 0.54
when the gold passage is prose; \m{nli\_set} barely changes (0.52 vs.\ 0.56). Computation is not the cause on FinanceBench: answers whose number
must be calculated are no harder for NLI (0.59). Part of the remaining ``error'' is ambiguity in the gold label: filings repeat figures across pages,
in 25\% of answers a distractor also contains the gold number, and 30\% of similarity failures pick such a passage.

\textbf{Controlled re-scoring} (Table~\ref{tab:why}). Rebuilding table rows with their column headers (``Depreciation and amortization: 2018 = 1,488;
2017 = 1,544; \dots'') raises \m{sim\_sent} by +0.103 [0.060, 0.147]; stripping boilerplate (running headers, page numbers) slightly hurts
($-$0.035), so boilerplate is not a cause. The largest gain comes from the numbers themselves. A \emph{number-aware} scorer ranks the passages by the
share of the sentence's numbers (years excluded, 1\% tolerance, unit shifts allowed) that they contain and lets an existing method break ties and
decide sentences without numbers. It needs no labels and no model. With similarity breaking ties it raises FinanceBench attribution from 0.35 to
0.58 (+0.239 [0.192, 0.288]); combined with the 20-label mixture it reaches 0.645 (+0.123 [0.084, 0.163] over the mixture). It acts only on sentences
with numbers, so it leaves the other four domains unchanged (all differences within $\pm$0.005, every CI contains zero). Per model, the macro average
rises from 0.74--0.81 to 0.76--0.83 and the worst domain from 0.51--0.55 to 0.59--0.71. In natural retrieval order, where the mixture's NLI weight
for finance does not carry over (NLI 0.44 vs.\ similarity 0.48), number matching with similarity is the best variant (worst domain 0.56--0.68 against
0.33--0.52 for the mixture). Phi-4 and Mistral-Small, which were added after the scorer was designed, show the same gains.

\begin{table}[!t]
\caption{FinanceBench Attribution Accuracy (Oracle, Seven Models, Same Answers): Passage Text Given to the Scorer and Number-Aware Scoring}
\label{tab:why}
\centering
\footnotesize
\setlength{\tabcolsep}{3pt}
\begin{tabular}{@{}lccc@{}}
\toprule
Scorer / passage text & FinBench & HotpotQA & SQuAD \\
\midrule
\m{sim\_sent}, original text & .351 & .888 & .879 \\
\m{sim\_sent}, boilerplate stripped & .315 & .888 & .879 \\
\m{sim\_sent}, table rows rebuilt & .410 & .888 & .879 \\
\m{sim\_sent}, rows + column headers & .446 & .888 & .879 \\
\midrule
\m{nli\_set} & .528 & .817 & .865 \\
20-label mixture & .528 & .888 & .881 \\
Number match + \m{sim\_sent} & .579 & .889 & .882 \\
Number match + mixture & \textbf{.645} & .889 & .882 \\
\bottomrule
\multicolumn{4}{@{}p{0.95\columnwidth}@{}}{\scriptsize Answers and contexts are not regenerated; only the scorer's view of the passages changes.
HotpotQA and SQuAD are prose controls. Exploratory on FinanceBench; confirmed on TAT-QA (T7, Table~\ref{tab:tatqa}).}
\end{tabular}
\end{table}

\subsection{Replication on a Second Financial Dataset (TAT-QA)}
\label{sec:tatqa}
Table~\ref{tab:tatqa} reports the seven generators on TAT-QA, with the outcome of every hypothesis registered before the data was generated.

\textbf{The collapse replicates (T1 supported).} With the gold passage guaranteed, TAT-QA attribution is below the same model's general-text level
for 7 of 7 models with both methods: \m{sim\_sent} 0.47--0.54 against 0.86--0.91 (gaps $-$0.33 to $-$0.40) and \m{nli\_set} 0.56--0.68 against
0.81--0.86. The drop is smaller than on FinanceBench (0.32--0.39), consistent with TAT-QA's short, clean excerpts, but it is present for every model.

\textbf{The method reversal and the mixture transfer (T4, T5 supported).} \m{nli\_set} beats \m{sim\_sent} by +0.116 [0.062, 0.171] (Holm $p<0.001$),
as on FinanceBench. The weight fitted on FinanceBench ($w=0$), applied unchanged, beats the default for unseen domains for all seven models
(0.56--0.68 vs.\ 0.50--0.59), and a weight picked from 20 TAT-QA labels lands at $w=0$ or $0.1$ for every model.

\textbf{Tables and computation (T2, T3 partly supported).} Similarity attribution is much worse when the answer comes from a table than from text
(0.39 vs.\ 0.74; $-$0.358 [$-$0.479, $-$0.233], Holm $p<0.001$), confirming the FinanceBench diagnosis on new data. For NLI the table effect is
smaller and does not survive correction ($-$0.126, Holm $p=0.058$). Computation affects the two methods the other way round: within table-sourced
questions, NLI attribution is much worse for computed than for quoted answers ($-$0.251 [$-$0.340, $-$0.160], Holm $p<0.001$), while similarity is
not affected ($-$0.026, n.s.). An entailment model can verify a quoted figure but not a figure that has to be calculated.

\textbf{The number-aware scorer is confirmed on held-out data (T7 supported).} Designed on FinanceBench (Section~\ref{sec:why}) and registered before
TAT-QA was run, number matching raises \m{sim\_sent} from 0.506 to 0.748 ($+$0.242 [0.197, 0.290]) and \m{nli\_set} from 0.622 to 0.759 ($+$0.137
[0.110, 0.166]), both Holm $p<0.001$, for every model (Table~\ref{tab:tatqa}). It closes most of the table gap: with it, table-sourced answers reach
0.71--0.74 and text-sourced ones 0.74--0.75.

\textbf{Refusal (T6 not supported at $\alpha=0.05$).} With the gold document removed, the share of refusals rises from 0.05 (1.5B) to 0.81 (32B)
(Spearman $\rho=0.67$, one-sided $p=0.051$), so the scale trend is suggestive but misses the registered threshold; Phi-4 and Gemma-27B refuse less
often (0.45--0.48) than models of similar size. Over-refusal does not fall with size ($\rho=+0.54$), as in the main study.

\begin{table*}[!t]
\caption{TAT-QA (200 Questions, Gold Passage Guaranteed): Accuracy, Attribution Accuracy With and Without Number Matching, and Refusal When the Gold
Document Is Removed (Trust Prompt)}
\label{tab:tatqa}
\centering
\footnotesize
\setlength{\tabcolsep}{5pt}
\begin{tabular}{@{}lccccccc@{}}
\toprule
& & \multicolumn{4}{c}{Attribution accuracy} & \multicolumn{2}{c}{Refusal} \\
\cmidrule(lr){3-6}\cmidrule(l){7-8}
Model & Accuracy & \m{sim\_sent} & \m{nli\_set} & number + \m{sim\_sent} & number + \m{nli\_set} & evidence removed & over-refusal \\
\midrule
Qwen-1.5B & .450 & .501 & .613 & .721 & .761 & .045 & .000 \\
Llama-8B & .565 & .503 & .568 & .758 & .759 & .580 & .120 \\
Qwen-14B & .595 & .521 & .684 & .817 & .815 & .590 & .090 \\
Phi-4-14B & .645 & .537 & .601 & .711 & .726 & .450 & .035 \\
Mistral-24B & .605 & .494 & .682 & .759 & .774 & .675 & .190 \\
Gemma-27B & .670 & .473 & .558 & .686 & .686 & .475 & .045 \\
Qwen-32B & .660 & .509 & .645 & .783 & .790 & .810 & .175 \\
\midrule
General text (same models, oracle) & -- & .86--.91 & .81--.86 & -- & -- & -- & -- \\
\bottomrule
\multicolumn{8}{@{}p{0.95\textwidth}@{}}{\scriptsize Pre-registered outcomes: T1 (TAT-QA below general text, 7/7 models, both methods) supported;
T2 (table $<$ text) supported for \m{sim\_sent}, not for \m{nli\_set} after Holm; T3 (computed $<$ quoted) supported for \m{nli\_set}, not for
\m{sim\_sent}; T4 (\m{nli\_set} $>$ \m{sim\_sent}) supported; T5 (FinanceBench weight transfers) supported; T6 (refusal rises with size) not
supported, $p=0.051$; T7 (number matching helps) supported for both methods. Over-refusal = refusals with the gold passage present (\m{oracle} context).}
\end{tabular}
\end{table*}

\subsection{LLM Judges as Attributors}
\label{sec:judge}
An obvious alternative to scoring passages is to ask an LLM which passage supports each sentence. We gave the same 12,008 oracle-condition sentences
(all seven generators, six datasets) to two judges, Llama-3.1-8B and Qwen2.5-32B, with a fixed prompt (``Which passage supports the sentence?
Reply with the passage number only''), at temperature 0. The plan (L1--L4) was registered before any judge output existed. One amendment was made
after inspecting reply formats but before computing any accuracy: Llama writes ``[3]'', so its first output token carries no passage number, and the
pick is therefore read from the written reply rather than from the first token's probabilities. The literal rule changes the judges' accuracies by
up to 0.05 (PubMedQA) and none of the verdicts below; under it the 32B judge leads the 8B judge by 0.041 [0.028, 0.055], significant but below our
0.05 threshold.

Table~\ref{tab:judge} shows that LLM judges are the most accurate attributors and do not collapse on financial text (0.69--0.79). The 32B judge beats
sentence similarity on FinanceBench (+0.340), TAT-QA (+0.284), CUAD (+0.094) and HotpotQA (+0.055), all Holm $p<0.01$, but not on SQuAD (+0.027) or
PubMedQA ($-$0.029), where it declines to name a passage for 16\% of sentences (often a bare ``Maybe''); L1 holds in 4 of 6 domains. A 4$\times$
larger judge is not better (L4: +0.003 [$-$0.010, 0.015]). The free number-aware scorer is \emph{not} equivalent to a judge (L2 fails): the 32B judge
leads it by 0.113 [0.062, 0.162] on FinanceBench and 0.041 [0.020, 0.062] on TAT-QA. It nevertheless closes 67--86\% of the gap between similarity
and the judge on FinanceBench and 85--87\% on TAT-QA, at no GPU cost, while a judge needs about 1 (8B) to 4 (32B) GPU-seconds per sentence on our
hardware. Adding number matching to a judge does not help meaningfully (L3: +0.027 and +0.008).

\begin{table}[!t]
\caption{Attribution Accuracy of Cheap Scorers and LLM Judges (Oracle Condition, 12,008 Sentences From Seven Generators)}
\label{tab:judge}
\centering
\footnotesize
\setlength{\tabcolsep}{2.5pt}
\begin{tabular}{@{}lcccccc@{}}
\toprule
Scorer & FinBench & TAT-QA & CUAD & PubMed & Hotpot & SQuAD \\
\midrule
\m{sim\_sent} & .352 & .506 & .677 & .820 & .888 & .881 \\
20-label mixture & .528 & .560 & .701 & .820 & .888 & .881 \\
Number + \m{sim\_sent} & .579 & .748 & .680 & .815 & .889 & .882 \\
Number + mixture & .645 & .753 & .701 & .815 & .889 & .882 \\
\midrule
Llama-3.1-8B judge & \textbf{.731} & .762 & .701 & \textbf{.846} & .940 & .908 \\
Qwen2.5-32B judge & .692 & \textbf{.789} & \textbf{.771} & .791 & \textbf{.943} & .908 \\
\bottomrule
\multicolumn{7}{@{}p{0.95\columnwidth}@{}}{\scriptsize Cost: cheap scorers need no GPU (\m{sim\_sent}: a small embedding model); judges about 1 (8B) and
4 (32B) GPU-seconds per sentence on 2$\times$T4. The 32B judge names no passage for 6.7\% of sentences (counted as misses).}
\end{tabular}
\end{table}

\subsection{A Second Medical Dataset With Sentence-Level Evidence (BioASQ)}
\label{sec:bioasq}
With the gold passage guaranteed, sentence similarity reaches 0.62--0.71 on BioASQ against 0.69--0.89 on PubMedQA for the same models: lower for 7 of
7 (B1 supported; Table~\ref{tab:bioasq}). PubMedQA's whole-abstract gold therefore made medical attribution look almost as easy as general text; with
snippet-level evidence, medicine sits between general text (0.83--0.93) and finance (0.32--0.39), like legal text. There is no finance-style failure:
similarity and NLI are about equal (B2: +0.023, n.s.), and number matching changes nothing (B3: +0.004, inside $\pm$0.05, supported), as it should on
prose. Retrieval costs less than in law and finance (B4: +0.037, not supported; mostly on factoid questions, whose gold is retrieved 80\% of the time).
Refusal transfers to this domain as well (B5 supported): every model of 8B or more reaches a balanced accuracy of 0.75--0.81 with low over-refusal
(0.02--0.07), while the 1.5B model stays near chance (0.52).

\begin{table}[!t]
\caption{BioASQ (200 Questions): Oracle Attribution With Sentence-Level Gold vs.\ PubMedQA's Abstract-Level Gold, and Refusal When the Evidence Is Removed}
\label{tab:bioasq}
\centering
\footnotesize
\setlength{\tabcolsep}{3pt}
\begin{tabular}{@{}lccccc@{}}
\toprule
& \multicolumn{2}{c}{BioASQ} & PubMedQA & \multicolumn{2}{c}{Refusal (trust prompt)} \\
\cmidrule(lr){2-3}\cmidrule(lr){4-4}\cmidrule(l){5-6}
Model & \m{sim\_sent} & \m{nli\_set} & \m{sim\_sent} & bal.\ acc. & over-ref. \\
\midrule
Qwen-1.5B & .713 & .682 & .883 & .524 & .006 \\
Llama-8B & .636 & .635 & .741 & .799 & .066 \\
Qwen-14B & .700 & .671 & .893 & .749 & .048 \\
Phi-4-14B & .662 & .635 & .795 & .778 & .048 \\
Mistral-24B & .668 & .639 & .803 & .805 & .018 \\
Gemma-27B & .624 & .593 & .685 & .793 & .018 \\
Qwen-32B & .674 & .659 & .886 & .775 & .054 \\
\bottomrule
\multicolumn{6}{@{}p{0.95\columnwidth}@{}}{\scriptsize Refusal: 167 question pairs per model (gold retrieved in the natural context vs.\ all listed
documents removed). PubMedQA values from Table~\ref{tab:attr}.}
\end{tabular}
\end{table}

\subsection{Evidence Removal and Refusal Across Model Scale}
\label{sec:refusal}
The removal condition gives unanswerable questions in every domain. Table~\ref{tab:refusal} reports the balanced accuracy of the model's own refusal
decision with the trust prompt (removed question vs.\ the same question with its evidence present). The 1.5B model barely departs from never refusing
(0.52--0.60). From 8B onwards the ability is present and works in every domain: on the shared subset the mean rises from 0.54 (1.5B) to 0.76 (8B),
0.77 (both 14B models), 0.83 (Mistral-24B), 0.80 (Gemma-27B) and 0.81 (32B). Within the Qwen family, balanced accuracy rises from 1.5B to 14B in 5 of 5
domains and from 14B to 32B in 4 of 5, which supports the first part of H4.

The second part of H4 is not supported. Larger models do not refuse less often when the evidence is present: Qwen-32B declines 32--40\% of answerable
HotpotQA, CUAD and FinanceBench questions whose evidence was in the context, more than the 8B and 14B models, and Mistral-24B declines 32\% in
FinanceBench. Gemma-27B and Phi-4 over-refuse least (0.08 and 0.11 on average). Refusal behaviour is therefore not just a function of size; it also
depends on the model family. Without the refusal instruction (vanilla prompt), all models answer most questions whose evidence was removed. The same
pattern appears on TAT-QA (Section~\ref{sec:tatqa}).

\begin{table}[!t]
\caption{Balanced Accuracy of Prompted Refusal on the Evidence-Removal Task (Trust Prompt; 0.5 = Never Refuses) and Over-Refusal on the Shared Subset}
\label{tab:refusal}
\centering
\footnotesize
\setlength{\tabcolsep}{2.5pt}
\begin{tabular}{@{}lcccccc@{}}
\toprule
Model & SQuAD & Hotpot & PubMed & CUAD & FinBench & Mean$^\ddagger$ \\
\midrule
\multicolumn{7}{@{}l}{\emph{Balanced accuracy}} \\
Qwen-1.5B & .521 & .560 & .530 & .598 & .518 & .544 \\
Llama-8B & .739 & .749 & .742 & .778 & .768 & .755 \\
Qwen-14B & .856 & .788 & .793 & .803 & .679 & .767 \\
Phi-4-14B$^\dagger$ & .750 & .786 & .874 & .829 & .605 & .769 \\
Mistral-24B$^\dagger$ & .786 & .812 & .929 & .842 & .763 & .827 \\
Gemma-27B$^\dagger$ & .774 & .828 & .889 & .868 & .632 & .798 \\
Qwen-32B$^\dagger$ & .845 & .786 & .879 & .776 & .737 & .805 \\
\midrule
\multicolumn{7}{@{}l}{\emph{Over-refusal (evidence present), shared subset}} \\
Qwen-1.5B & .000 & .042 & .010 & .000 & .000 & .010 \\
Llama-8B & .095 & .365 & .061 & .237 & .158 & .183 \\
Qwen-14B & .119 & .312 & .020 & .237 & .211 & .180 \\
Phi-4-14B$^\dagger$ & .024 & .281 & .010 & .079 & .158 & .110 \\
Mistral-24B$^\dagger$ & .000 & .177 & .030 & .211 & .316 & .147 \\
Gemma-27B$^\dagger$ & .000 & .167 & .051 & .026 & .158 & .080 \\
Qwen-32B$^\dagger$ & .119 & .375 & .030 & .395 & .316 & .247 \\
\bottomrule
\multicolumn{7}{@{}p{0.95\columnwidth}@{}}{\scriptsize Pairs with the gold chunk retrieved (natural order). Balanced accuracy for 1.5B--14B on all
questions; $^\dagger$ 100 questions per domain; $^\ddagger$ mean over domains on the shared subset. Pairs per domain on the subset: 42, 96, 99, 38, 19.}
\end{tabular}
\end{table}

\textbf{Missing evidence versus wrong answers.} Confidence signals separate ``evidence present'' from ``evidence removed'' much better than they
separate correct from wrong answers: for Llama-8B the best AUROC per domain is 0.77--0.95 for the first task and 0.64--0.76 for the second. A threshold
fitted on the two general datasets transfers to the specialised domains with a mean loss of 0.035--0.050 balanced accuracy. The risk-controlled setting
behaves differently. A refusal threshold calibrated on general questions to keep the answered-and-wrong rate below $\alpha=0.2$ \cite{angelopoulos2024crc}
kept the target on SQuAD (0.12--0.14) but broke it in every specialised domain for both small models (0.22--0.34), while 20 labelled in-domain questions
restored it (0.12--0.18). Detecting that evidence is missing transfers across domains; predicting whether an answer is correct does not.

\subsection{Confidence Signals and Token Log-Probabilities}
\label{sec:lp}
Table~\ref{tab:lp} summarises the pre-registered analysis over 70 cells (7 models $\times$ 2 systems $\times$ 5 domains). \m{lp\_mean} is significantly
above chance in 52 of 70 cells (the control in 58) and is higher than the control in 41. Its practical value is in selective answering: keeping only
the 50\% most confident answers raises accuracy by 6--14 points depending on the model, more than retrieval similarity (0.5--8.2 points) for every
model. Pooled over models the gain is +15.0 points in CUAD, +14.3 in FinanceBench, +12.8 in PubMedQA, +7.9 in HotpotQA and +3.6 in SQuAD.

The other hypotheses are not supported. \textbf{H1}: mean AUROC is 0.639 (1.5B), 0.674 (8B), 0.676 (Qwen-14B), 0.674 (Phi-4-14B), 0.667 (24B),
0.650 (27B) and 0.647 (32B). It rises from 1.5B to 8B and then stays flat or falls, in every family. The upward trend we saw with the two small models
alone was not a scale trend. \textbf{H2}: at $\lambda\geq2$ only 2 of 140 tests for \m{lp\_mean} survive the false-discovery control; no signal
reliably beats the trivial policies when wrong answers are costly. \textbf{H3} holds only in part: log-probabilities are better than the other signals
on HotpotQA and PubMedQA as predicted, but they are also better on CUAD and FinanceBench, where we had expected retrieval signals to match them.

\begin{table}[!t]
\caption{Token Log-Probabilities as a Confidence Signal: Mean AUROC for a Correct Answer and Accuracy Gain When Answering Only the 50\% Most Confident Questions}
\label{tab:lp}
\centering
\footnotesize
\setlength{\tabcolsep}{3pt}
\begin{tabular}{@{}lcccc@{}}
\toprule
& \multicolumn{2}{c}{AUROC (mean over cells)} & \multicolumn{2}{c}{Accuracy gain at 50\%} \\
\cmidrule(lr){2-3}\cmidrule(l){4-5}
Model & \m{lp\_mean} & control & \m{lp\_mean} & \m{ret\_top1} \\
\midrule
Qwen-1.5B & .639 & .656 & +.088 & +.049 \\
Llama-8B & .674 & .651 & +.087 & +.054 \\
Qwen-14B & .676 & .640 & +.116 & +.056 \\
Phi-4-14B$^\dagger$ & .674 & .647 & +.117 & +.035 \\
Mistral-24B$^\dagger$ & .667 & .672 & +.104 & +.082 \\
Gemma-27B$^\dagger$ & .650 & .646 & +.136 & +.039 \\
Qwen-32B$^\dagger$ & .647 & .630 & +.060 & +.005 \\
\bottomrule
\multicolumn{5}{@{}p{0.95\columnwidth}@{}}{\scriptsize AUROC over vanilla and improved-retrieval systems and five domains; ``control'' is the best
non-log-probability signal of each cell, chosen on the same data. Gains with the vanilla prompt. $^\dagger$ 100 questions per domain.}
\end{tabular}
\end{table}

\subsection{Hallucination Detection on RAGTruth and Transfer}
\label{sec:ragtruth}
Table~\ref{tab:ragtruth} reports response-level detection on the RAGTruth QA test split (900 responses) \cite{niu2024ragtruth}. Our zero-shot detector
(\m{nli\_set} with DeBERTa-v3-base, threshold tuned on the training split) reaches 43.6 F1. A 3$\times$ larger NLI model does not help (43.0). A
ModernBERT-base checker \cite{warner2024modernbert} with about 150M parameters, fine-tuned on the RAGTruth training split, reaches 51.4 F1 and an AUROC of
0.825. A paired bootstrap on the same 900 responses confirms the improvement over our zero-shot detector: +7.8 F1 [3.3, 12.2], $p=0.001$, and +0.065
AUROC [0.020, 0.109], $p=0.002$. The checker is above the published GPT-4-turbo prompting result and 17 points below a fine-tuned Llama-2-13B that has
about 90 times more parameters.

On our own domains the checker does not help. Its scores are saturated (mean probability of ``hallucinated'' 0.73--0.88 per system, against 0.24 on
RAGTruth), its AUROC for answer correctness is 0.46--0.63 and never the best signal in 30 of 30 cells, and it ranks last of 12 signals at detecting
removed evidence. Likely reasons are input length (about 480 tokens in RAGTruth QA against about 1,300 in our contexts), older generators in RAGTruth,
and span-level labels that differ from short-answer correctness in specialised text.

\begin{table}[!t]
\caption{Response-Level Hallucination Detection on the RAGTruth QA Test Split}
\label{tab:ragtruth}
\centering
\footnotesize
\begin{tabular}{@{}lcc@{}}
\toprule
Detector & F1 & AUROC \\
\midrule
GPT-3.5-turbo prompting \cite{niu2024ragtruth} & 30.8 & -- \\
SelfCheckGPT \cite{niu2024ragtruth} & 43.7 & -- \\
GPT-4-turbo prompting \cite{niu2024ragtruth} & 45.6 & -- \\
Fine-tuned Llama-2-13B \cite{niu2024ragtruth} & 68.2 & -- \\
\midrule
Ours: similarity, zero-shot & 35.2 & .620 \\
Ours: chunk NLI, zero-shot (base) & 40.3 & .720 \\
Ours: window NLI, zero-shot (base) & 42.8 & .738 \\
Ours: set NLI, zero-shot (base) & 43.6 & .760 \\
Ours: chunk NLI, zero-shot (large) & 43.0 & .762 \\
Ours: ModernBERT-base checker, supervised & \textbf{51.4} & \textbf{.825} \\
\bottomrule
\multicolumn{3}{@{}p{0.95\columnwidth}@{}}{\scriptsize Checker vs.\ set NLI (base), paired bootstrap: F1 +7.8 [3.3, 12.2], $p=0.001$.
Published rows: point estimates only.}
\end{tabular}
\end{table}

\subsection{Human Validation of the Automatic Labels}
\label{sec:human}
To check the automatic support and attribution labels, we sampled 100 answer sentences (20 per domain) from the vanilla and oracle runs of the 1.5B and
8B models. Each sentence was labelled \emph{Supported}, \emph{Partially supported} or \emph{Not supported}, together with its supporting passages and a
short quoted reason. The labels were proposed by an LLM (Claude), which is not one of the evaluated generators, and verified independently by two of the authors, who are not
domain experts. The verifiers agreed with 91.0\% and 82.0\% of the proposed labels (Cohen's kappa between them 0.369); disagreements were resolved by discussion.
Of the 100 sentences, 55 are supported, 29 partially supported and 16 not supported; 9 of the 16 unsupported sentences are from FinanceBench.
Table~\ref{tab:human} compares the methods with these labels. Entailment-based methods detect unsupported sentences better than similarity (AUROC
0.704 for \m{nli\_window} and 0.713 for \m{nli\_set}, against 0.673 for \m{sim\_sent}), which supports keeping NLI for verification. For
\emph{locating} the supporting passage, similarity is as good as entailment: the top-scored chunk is among the labelled supporting passages for 0.85--0.87
of supported sentences with \m{sim\_sent}, \m{sim\_chunk} and \m{nli\_set}, against 0.74 for \m{nli\_chunk}. Finally, the automatic gold chunk used to
compute attribution accuracy is among the labelled supporting passages for 94.0\% (47 of 50) of supported sentences whose context contained a gold
chunk, which supports attribution accuracy as a metric.

\begin{table}[!t]
\caption{Automatic Methods Against the Sentence-Level Support Labels (100 Sentences, 20 per Domain)}
\label{tab:human}
\centering
\footnotesize
\begin{tabular}{@{}lcc@{}}
\toprule
Method & AUROC (supported) & Top chunk in labelled set \\
\midrule
\m{sim\_chunk} & .610 & .870 \\
\m{sim\_sent} & .673 & .870 \\
\m{nli\_chunk} & .618 & .741 \\
\m{nli\_window} & \textbf{.704} & .815 \\
\m{nli\_set} & .713 & .852 \\
\midrule
Automatic gold chunk & -- & .940 \\
\bottomrule
\multicolumn{3}{@{}p{0.95\columnwidth}@{}}{\scriptsize AUROC for separating \emph{Supported} from the other two labels, using each
method's best score for the sentence. Last column: share of supported sentences whose top-scored chunk is one of the labelled supporting passages.}
\end{tabular}
\end{table}

\subsection{Negative Results}
\label{sec:neg}
\textbf{An end-to-end trust layer does not beat trivial policies.} Our first ``full'' system (trust prompt, an NLI hallucination filter and a confidence
gate fitted on SQuAD) seemed to improve a per-question utility by up to +0.73 over plain RAG. Under the protocol of Section~\ref{sec:protocol}, refusing
every question already scores 0.50 on SQuAD and 0.33 on CUAD, only 1 of 20 cells has a usable signal that beats the best trivial policy at $\lambda=1$,
and the full system scores below plain trust-prompt answers on HotpotQA, PubMedQA and CUAD. The apparent gains came from refusing many questions in
domains that contain many unanswerable ones.

\textbf{An order artefact in evidence-set NLI.} Giving the joint score to both chunks of a pair and breaking ties toward the lower index made
\m{nli\_set} depend on where the gold chunk appeared (0.152 difference between positions 1--2 and 4--5, against about 0.04 for the other methods).
Breaking ties by each chunk's own score removes the effect. With the fix, \m{nli\_set} still beats single-chunk window NLI (0.632 vs.\ 0.579) but not
sentence similarity (0.744) on average for attribution.

\textbf{A withdrawn claim and a rejected hypothesis.} An exploratory sweep that counted ``wins'' of 16 signals per cell without multiplicity control
suggested that log-probabilities help at high error cost; the pre-registered test does not support it. The two-model trend that suggested
log-probabilities improve with model size (H1) did not hold with seven models. On TAT-QA, the registered scale trend in refusal (T6) narrowly missed
its threshold ($p=0.051$), and the table effect on NLI attribution (T2) did not survive correction; both are reported as not supported.

% =====================================================================================================================
\section{Discussion and Limitations}
\label{sec:limits}
\textbf{What travels and what does not.} Detecting \emph{missing} evidence travels across domains: once a model is large enough, its prompted refusal
works similarly in all five domains and on BioASQ, and thresholds on retrieval or NLI signals transfer with small losses. Predicting \emph{correctness}
does not travel: risk-controlled thresholds, a supervised detector trained on general data and our earlier end-to-end system all fail on specialised
text. Attribution lies in between, and how far it travels depends on the scorer. Cheap scorers work on general text, are moderate on medical and
legal text (once medical evidence is measured at sentence level) and fail on financial documents for every model, on two financial datasets. The
failure has a concrete cause: embedding similarity cannot match an answer sentence to table cells, and entailment cannot verify a computed figure.
Matching the numbers repairs much of it at no cost; LLM judges do not have the problem at all, at about 1--4 GPU-seconds per sentence. 20 labelled
questions per domain are enough to choose between similarity and entailment for the rest.

\textbf{What model size does and does not change.} Larger models refuse better when evidence is missing and, mostly between 1.5B and 8B, answer more questions correctly when
the evidence is guaranteed. They do not fix retrieval failures, they do not close the attribution gap in finance (neither does changing the model
family), their log-probabilities are not more informative, and some of them over-refuse more.

\textbf{Practical advice.} For a team building a RAG system for a specialised domain, our results suggest: invest first in domain-specific retrieval;
for attribution, use an 8B LLM judge if its cost is acceptable (a 32B judge was not better); otherwise match the numbers in each sentence first, then
use cheap sentence similarity and check 20 labelled questions to decide whether entailment is needed; feed tables to any scorer as rows with their
headers rather than as raw PDF text; measure medical attribution against sentence-level evidence, not whole abstracts; use a model of at least about 8B
parameters if it must refuse on its own, and measure over-refusal; and use token log-probabilities to answer only the most confident questions rather
than as a guarantee of correctness.

\textbf{Limitations.} (i) The largest model has 32B parameters, and four of the seven models answer 100 questions per domain. (ii) Labels are automatic.
They are validated on 100 sentences whose labels were proposed by an LLM and checked by non-expert authors, not by independent expert annotators;
checking a proposal can make agreement look higher than independent labelling would. (iii) PubMedQA's whole-abstract gold makes its attribution and
removal results easy; BioASQ (200 questions) corrects the attribution picture but is a single, smaller sample. (iv) Some analyses rest
on few questions: FinanceBench refusal pairs (19--28), and leave-one-out (40 answers per domain for two models). (v) Each configuration was run once,
apart from the noise-floor check, and the supervised checker was trained once. (vi) All models were run with 4-bit weights. (vii) The number-aware
scorer and the table-row rebuilding were designed on FinanceBench; they are confirmed on one held-out dataset (TAT-QA, 200 questions, 16 of them
text-computed), and TAT-QA uses only the oracle and removal conditions because its questions do not name the company. (viii) The row rebuilder is a
heuristic for one-cell-per-line PDF text, and the effect of rebuilt rows on NLI scoring was not measured. (ix) The LLM judges use one fixed prompt
with forced choice; the 32B judge's refusals to name a passage (6.7\%) count as misses, and the reading rule was amended once, before accuracies were
computed (Section~\ref{sec:judge}).

% =====================================================================================================================
\section{Conclusion}
We evaluated attribution, hallucination detection and refusal for RAG under one fixed pipeline across general, medical, legal and financial documents,
using gold evidence, counterfactual contexts and seven generators from five families and 1.5B to 32B parameters. Retrieval is the largest loss in
legal and financial text. Even with retrieval taken out, attribution fails on financial documents for every model, on two datasets, and the best
method changes with the domain. The cause is the tables and numbers: a label-free number-aware scorer, designed on FinanceBench and confirmed on
held-out TAT-QA data under pre-registered tests, raises financial attribution by 0.12--0.24 without affecting other domains and closes most of the
gap to LLM judges at no cost, and 20 labelled questions per domain are enough to choose the rest of the scorer. LLM judges remain the most accurate
attributors, and an 8B judge is as good as a 32B one. Measured against sentence-level evidence (BioASQ), medical attribution is clearly harder than
whole-abstract gold suggests. Refusing when evidence is missing improves with model
size and then works in every domain, but over-refusal does not go away. Predicting whether an answer is correct stays hard everywhere: token
log-probabilities help to answer only the most confident questions but do not improve with size, and a hallucination checker that clearly improves on a
public benchmark does not transfer. The code, run configurations, pre-registered analysis plans and aggregate results are available at
\url{https://github.com/Surveerssg/trustworthy-rag-attribution}; all datasets are public (BioASQ after registration), and the released scripts
rebuild the gold-evidence labels from them.

\section*{Acknowledgment}
All experiments were run on free Kaggle GPU resources. The authors thank the Department of Computer Science and Engineering, NSUT, for its support.

\bibliographystyle{IEEEtran}
\bibliography{references}

\end{document}
